% =====================================================================
% Capítulo 4. Etapa 1: bit flipping en LoRaWAN 1.1  Presupuesto: 9 páginas
% ---------------------------------------------------------------------
% Este capítulo tiene método Y resultados de la etapa 1 (es una etapa cerrada).
% Su función en la historia: demostrar que (1) un atacante sin la clave puede
% causar daño real con valores plausibles, (2) la medición aislada no basta y
% (3) el contexto temporal-espacial ayuda. Eso motiva la etapa 2.
% Fuentes en el repo: docs/01_umbrales_normativos.md, docs/02_codificacion.md,
% docs/protocolo_lstm_v4_multianual.md, results/lstm_v4_multianual_v1/,
% experiments/cayenne_mod/cierre_distribucion_pesos.md.
% =====================================================================
\chapter{Etapa 1: daño y detectabilidad del \textit{bit flipping}}
\label{cap:bit-flipping}

\section{Modelo de amenaza}
\label{sec:amenaza}
% Debe contestar: quién es el atacante, dónde está (tramo NS -> AS), qué sabe
% (estructura CayenneLPP), qué NO sabe (claves), qué quiere (cambiar la banda
% comunicada). Figura: figuras/modelo_amenaza.tex.

\section{Generación sistemática de ataques}
\label{sec:generacion}
% Debe contestar: cómo se codifica, cifra, modifica y descifra; qué posiciones
% de bit se recorren; qué se guarda de cada muestra (valor real, recibido, bit,
% dirección, daño). Debe permitir que otro lo reproduzca.

\section{Definición del daño}
\label{sec:dano}
% Ecuación D = 1[f(x_real) != f(x_recibido)] con f = banda NOM-172.
% Ejemplo: 32 ppb lejos del corte no dañan; 2 ppb junto al corte sí.

\section{Detección: medición aislada contra contexto}
\label{sec:deteccion-e1}
% Capa 1 (medición aislada) y LSTM con contexto espacial; detector V4
% (31 LSTM, ajuste 2020--2022, calibración 2023, prueba 2024--2026).

\section{Resultados}
\label{sec:resultados-e1}
% - Tabla: detección con 32 y 64 ppb, aislada contra LSTM (estación CCA, 2025).
% - Óptimo interior del atacante: máximo daño no detectado en medio.
% - Bits 3--5 concentran el daño no detectado.
% - Deriva estacional -> perfil causal de 14 días (se hereda a la etapa 2).
% - Contramedida en CayenneLPP: resultado negativo (255 -> 567 cambios de
%   banda no detectados en 179 399 intentos).

\section{De la etapa 1 a la etapa 2}
\label{sec:puente}
% Media página: qué se aprendió y por qué el problema se generaliza a
% «qué estación reporta datos anómalos», sin importar la causa.
